\documentclass{article}

% Shared preamble from project root
\usepackage{lecture}

\begin{document}

% \begingroup
%     \hypersetup{linkbordercolor=white, linkcolor=black} 
%     \tableofcontents
%     \newpage
% \endgroup

\section{Thread Safety}\label{sec:ts_thread_safety} % (fold)
\begin{definition}{Thread Safe}{tsafe}
    A class is \textbf{thread safe} if continues to behave correctly when accessed from multiple threads, regardless of the scheduling or interleaving of the execution of those threads by the runtime environment.
    \begin{note}
        Without additional synchronization or other coordination on the part of the calling code
    \end{note}
\end{definition}
Correctness means that a class conforms to its specification, a good specification defines:
\begin{itemize}
    \item \textbf{invariants} constraining an object's state;
    \item \textbf{post-conditions} describing the effects if its operations;
\end{itemize}
\begin{note}
    Thread-safe classes encapsulate any needed synchronization, so the clients don't need to provide their own
\end{note}
%
\subsection{Shared Mutable State}\label{sub:ts_shared_mutable_state} % (fold)
Writing thread-safe code is -- at its core -- about \textit{managing access to state} and in particular to \textbf{shared} and \textbf{mutable state}:
\begin{definition}{Shared}{shared}
    A variable or object could be accessed by multiple threads.
\end{definition}
\begin{definition}{Mutable}{mutable}
    Its value could change during its lifetime.
\end{definition}
\begin{critical}[Race Conditions]
    If multiple threads access the same \hyperref[def:mutable]{Mutable} state variable without appropriate synchronization, the program is broken
\end{critical}
There are three ways to fix it:
\begin{itemize}
    \item don't save state variable across threads;
    \item make the state variable across threads;
    \item use synchronization whenever accessing the state variable;
\end{itemize}
% subsection Shared Mutable State (end)
%
\subsubsection{Stateless \& Stateful Object}\label{sub:stateless_stateful_object} % (fold)
\textbf{Stateless} objects are always \hyperref[def:tsafe]{Thread Safe}: actions of a thread accessing a stateless object cannot affect the correctness of operations in other threads.
\begin{note}
    Thread safety is undermined as soon as we share and access in R/W \textbf{stateful} objects: \hyperref[def:mutable]{Mutable} state-full objects in particular
\end{note}
% subsubsection Stateless \& Stateful Object (end)
%
\subsection{Race Conditions}\label{sub:race_conditions} % (fold)
The concurrent execution of non-atomic sequence of statements that should be considered atomic, generate \textbf{race condition}.
\begin{note}
    It occurs when the \textbf{correctness of a computation} depends on the relative timing or interleaving of multiple threads by the runtime
\end{note}
%
\subsubsection{Lost Updates \& Check-and-Act}\label{sub:lost_updates_check_and_act} % (fold)
\begin{definition}{Lost Updates}{lost-update}
    When executing concurrently non-atomic read-modify-write operations (\texttt{count++}).
\end{definition}
\begin{definition}{Check-and-Act}{check-act}
    When a potentially state observation is used to make a decision on what to do next.
    \begin{note}
        Since check and act are not atomic, the state can change after check and before act
    \end{note}
    \begin{example}{}{}
        \begin{codeblock}[txt]{}
        if (file X doesn't exists)  -- check
            then create file X      -- act
        \end{codeblock}
    \end{example}
\end{definition}
\begin{note}[Compound Atomic Actions]
    \hyperref[def:check-act]{Check-and-Act} and read-modify-write are examples of compound actions: the sequence of operation that must be executed atomically in order to remain \hyperref[def:tsafe]{Thread Safe}
\end{note}
% subsubsection Lost Updates \& Check-and-Act (end)
% subsection Race Conditions (end)
% section Thread Safety (end)
\end{document}

